\scanpage{7}
\mspage{5}
schéma fini étale sur \(S\), alors on a des isomorphismes
\[
\operatorname{Hom}_{S}(\bar s,S')\xleftarrow{\sim}\operatorname{Hom}_S(\deleted{\operatorname{Spec}}\widetilde S,S')\xrightarrow{\sim}\operatorname{Hom}_S(\deleted{\operatorname{Spec}}\bar\eta,S')
\]
d'où un isom canonique
\[
\operatorname{Hom}_S(\bar s,S')\simeq\operatorname{Hom}_S(\bar\eta,S'),\tag*{(2.3.1)}
\]
fonctoriel en \(S'\) (isom. des \inserted{2} ``foncteurs fibres'' sur la cat. des \unsure{ouv.|rev.|var.} étales de \(S\), définis resp. par les pts géom. \(\unsure{\bar t|\bar\eta|\bar s}\) et \(\bar s\)). Faisant \(S'=\mu_n\), \(n\) \inserted{entier} inversible ds \(k\), on trouve en particulier un isom
\[
\mu_n(\bar s)={}_n\bar k^*\simeq\mu_n(\overline K)={}_n\overline K^*,\tag*{(2.\unsure{3|1}.2)}
\]
d'où par passage à la limite sur les puissances un nb premier \(\ell\), \(\ell\ne\operatorname{car.}k\) :
\[
T_\ell(\bar k^*)\simeq T_\ell(\overline K^*).\tag*{(2.3.3)}
\]
\deleted{On \unsure{notera|note}.} Les deux membres, \deleted{qui} sont des modules libres de rang 1 sur \(\mathbb Z_\ell\), \deleted{\unsure{les ident}} \inserted{qu'on identifie} désormais et désignera par \(\mathbb Z_\ell(1)\). \deleted{\unsure{Ses puissances}}
\[
\mathbb Z_\ell(1)\simeq\ldots\ldots\tag*{(2.3.3)}
\]
\inserted{Sa puiss.} \unsure{tensorielle|tensorielles} \(n\).ième (\(n\in\mathbb Z\)) est notée \(\mathbb Z_\ell(n)\). Donc on a
\[
\mathbb Z_\ell(n)\simeq\varprojlim_\nu\mu_{\ell^\nu}(\bar k^*)^{\otimes n}\simeq\varprojlim_\nu\mu_{\ell^\nu}(\overline K^*)^{\otimes n}.\tag*{(2.3.4)}
\]
Ces isomorphismes \deleted{\unsure{sont compatibles avec}} transforment les opérations de \(\pi\) sur \(\mu_n(\overline K^*)\) \inserted{etc} en celles de \(\pi_0\) sur \(\mu_n(\bar k^*)\) etc\ldots, compte tenu de l'épimorphisme canonique \(\pi\mathrel{\unsure{-|\to}}\pi_0\).
% archive sheet number: 6
